\chapter{Tests, validation et resultats}
\label{chap:tests}

\section{Strategie de test}

Le projet integre les quatre dependances de test standard de l'ecosysteme Spring
Boot couvrant respectivement la persistance, le web MVC, la securite et le
traitement par lots (\texttt{data-jpa}, \texttt{webmvc}, \texttt{security},
\texttt{batch}), completees par AssertJ et Hamcrest pour l'expressivite des
assertions. La strategie
de test suit la structure en couches du \cref{fig:layers} :

\begin{itemize}
  \item \textbf{Tests unitaires} sur les services de calcul (\texttt{Risk\-Calculation\-Service},
  \texttt{Dashboard\-Metrics\-Service}) : verification des formules du
  \cref{chap:modelisation} sur des jeux de donnees synthetiques controles, en
  particulier des cas limites (hote sans vulnerabilite, actif sans correspondance
  CMDB, ARS proche du seuil de saturation).
  \item \textbf{Tests d'integration Spring Batch}, permis par la dependance
  \texttt{batch-test}, pour valider que le pipeline
  d'ingestion (lecture StAX $\to$ mapping $\to$ ecriture JPA) produit bien les
  entites attendues a partir d'un extrait XML de reference.
  \item \textbf{Tests de mapping} (MapStruct) verifiant que chaque champ du DTO issu
  du parsing XML est correctement reporte sur l'entite JPA correspondante.
\end{itemize}

\section{Validation numerique du modele ARS}
\label{sec:validation-numerique}

Pour illustrer et verifier concretement les proprietes demontrees au
\cref{chap:modelisation}, cette section deroule le calcul complet sur un actif
fictif mais realiste.

\begin{exemple}[Calcul complet de l'ARS pour un serveur de base de donnees de production]
\label{ex:calcul-complet}
Soit un hote $h$ correspondant, dans la CMDB, a un CI de role
\texttt{DATABASE\_SERVER} en environnement \texttt{PRODUCTION}, supervise par le SOC,
sans EDR installe, et exclu du cycle de patch standard (maintenance fenetre
critique). Qualys a detecte sur cet hote $2$ vulnerabilites Critiques (QDS moyen
$92$), $3$ vulnerabilites Elevees (QDS moyen $78$) et $5$ vulnerabilites Moyennes
(QDS moyen $55$), et l'une des vulnerabilites correspond au QID $45004$ (exposition
Internet).

\smallskip
\noindent\textbf{1. ACS.} Avec des poids configures $\rho(\text{DATABASE\_SERVER})=0{,}9$
et $\varepsilon(\text{PRODUCTION})=1{,}0$ :
\[
ACS(h) = 1+4\cdot\frac{0{,}9+1{,}0}{2} = 1+4\times0{,}95 = 4{,}8 .
\]

\smallskip
\noindent\textbf{2. WeightedQDS} (\cref{eq:qds-impl}, poids de reference) :
\[
\begin{aligned}
\text{terme}_c &= 1{,}0 \times 92 \times 2^{1/100} = 92 \times 1{,}00696 \approx 92{,}64 \\
\text{terme}_h &= 0{,}2 \times 78 \times 3^{1/100} = 15{,}6 \times 1{,}01108 \approx 15{,}77 \\
\text{terme}_m &= 0{,}05 \times 55 \times 5^{1/100} = 2{,}75 \times 1{,}01630 \approx 2{,}79 \\
\text{terme}_l &= 0 \quad (\text{aucune vulnerabilite Faible}) \\
\mathrm{WeightedQDS}(h) &\approx 92{,}64+15{,}77+2{,}79 = 111{,}20
\end{aligned}
\]

\noindent Conformement a la \cref{prop:ratio-severite}, on verifie que le terme
Critique domine largement : il represente a lui seul
$92{,}64/111{,}20\approx 83\%$ du score, bien que ne comptant que $2$
vulnerabilites sur $10$.

\smallskip
\noindent\textbf{3. Facteurs de modulation} (\cref{tab:multiplicateurs}) : exposition
Internet averee, $I_{External}=1{,}2$ ; SOC present, $M_{SOC}=0{,}8$ ; pas d'EDR,
$M_{AV}=1{,}0$ ; exclu du cycle de patch, $M_{Patch}=1{,}3$ :
\[
M_{Controls}(h) = 0{,}8\times1{,}0\times1{,}3 = 1{,}04 .
\]

\smallskip
\noindent\textbf{4. ARS} (\cref{eq:ars}) :
\[
ARS(h) = \min\big(4{,}8 \times 111{,}20 \times 1{,}2 \times 1{,}04\;,\;1000\big)
= \min(666{,}5\,;\,1000) = 666{,}5 .
\]

\noindent Comme prevu par la \cref{prop:saturation} (seuil de saturation
$\approx128{,}2$ pour $\mathrm{WeightedQDS}$, tres inferieur au $111{,}20$ obtenu
ici mais du meme ordre de grandeur), l'actif s'approche de la saturation sans
l'atteindre : le plafond de $1000$ n'intervient pas, $ARS(h)$ est exactement le
produit calcule.

\smallskip
\noindent\textbf{5. Classification} (\cref{eq:classify}, \cref{tab:sla}) :
$600\le 666{,}5<700$, donc $h$ est classe \textbf{RED}, avec un SLA exige de
\textbf{7 jours}.

\smallskip
\noindent\textbf{6. Statut SLA.} Si la vulnerabilite Critique la plus ancienne a ete
detectee il y a $10$ jours et que le CI n'est pas en exception, alors la date
d'echeance ($10-7=3$ jours dans le passe par rapport a aujourd'hui, en supposant la
detection il y a $10$ jours et un SLA de $7$ jours) est depassee :
$\mathrm{StatutSLA}(h) = \text{VIOLATED}$ (\cref{eq:sla-status}). Sur le
\cref{fig:chart-sla}, cet actif se situerait au point $(666{,}5\,;\,10)$, au-dessus
de la marche « $7$ jours » de la zone RED : visuellement en violation, coherent avec
le calcul.
\end{exemple}

\section{Validation du MTTR pondere}

\begin{exemple}
Sur une periode donnee, les vulnerabilites corrigees (hors actifs en exception) se
repartissent ainsi : classe RED, $C_{red}=4$ corrections, $TTR$ moyen
$\mathrm{MTTR}_{red}=5$ jours ; classe ORANGE, $C_{orange}=10$ corrections,
$\mathrm{MTTR}_{orange}=12$ jours. Avec les poids de performance de reference
$w_{red}=5$ et $w_{orange}=3$ :
\begin{align*}
\mathrm{MTTR}_w &= \frac{w_{red}\,C_{red}\,\mathrm{MTTR}_{red} + w_{orange}\,C_{orange}\,\mathrm{MTTR}_{orange}}
                       {w_{red}\,C_{red} + w_{orange}\,C_{orange}}
= \frac{5\times4\times5 + 3\times10\times12}{5\times4+3\times10} \\
&= \frac{100+360}{20+30} = \frac{460}{50} = 9{,}2 \text{ jours}.
\end{align*}
Conformement a la \cref{prop:mttr-weighted-avg}, $\mathrm{MTTR}_w=9{,}2$ est bien
compris entre $\min(5,12)=5$ et $\max(5,12)=12$ ; il est de plus plus proche de la
valeur ORANGE que de la moyenne simple non ponderee
$\tfrac{5+12}{2}=8{,}5$ ne l'aurait laisse supposer, car le poids relatif effectif
de la classe ORANGE ($w_{orange}C_{orange}=30$) depasse celui de RED
($w_{red}C_{red}=20$) malgre son poids de performance $w$ individuel plus faible --
consequence directe du volume corrige plus important sur cette classe.
\end{exemple}

\section{Bilan des resultats obtenus}

La plateforme realisee couvre de bout en bout le scenario cible : import d'un
export Qualys et d'un export CMDB, calcul automatique et parallelise de l'ARS pour
chaque actif reconcilie, classification et gouvernance SLA, historisation des
instantanes de risque, restitution temps reel sous forme de cinq graphiques
complementaires (\cref{chap:visualisation}), et un prototype d'interrogation en
langage naturel via l'assistant RAG (\cref{chap:realisation}), ce dernier restant au
stade du developpement et de la conception et n'entrant pas dans le perimetre de la
campagne de validation ci-dessus. Les proprietes mathematiques demontrees au
\cref{chap:modelisation} (bornage de l'ARS, monotonie, coherence de la
classification, comportement en moyenne ponderee du MTTR) ont ete verifiees
numeriquement sur les exemples ci-dessus, confirmant que l'implementation respecte
fidelement le modele theorique specifie -- a l'exception assumee et documentee du
facteur $\mathrm{Global\_I_{Ops}}$ de l'ERS (\cref{sec:ecart-ers}).
